\section{Related Work}
\label{sec:related_work}

\textbf{Explicit memory for LLMs.} RAG systems \citep{lewis2020retrieval} retrieve documents into the context window, subject to retrieval quality and context dilution. MemGPT \citep{packer2023memgpt} manages a tiered memory buffer with explicit read/write operations. Reflexion \citep{shinn2023reflexion} accumulates verbal reflections as prompt prefixes, enabling cross-episode learning but with linearly growing context cost. Long-context approaches extend the window but do not change the fundamental architecture: memory remains content in context, subject to dilution and attention competition. All existing approaches operate as explicit memory --- the model reads stored content and can reason over it. HEXIS provides a complementary, \emph{non-context-mediated memory channel}: enmeshed modulation shapes attention and value extraction before the host model's chain-of-thought processes the content, so the modulation cannot be argued against by an adversary in the same conversation.

\textbf{Parameter-efficient adaptation.} LoRA \citep{hu2021lora} and its variants apply low-rank weight perturbations via gradient descent, producing fixed adapters per training run. Adapters \citep{houlsby2019parameter} insert bottleneck layers between frozen blocks. Prefix tuning \citep{li2021prefix} and prompt tuning \citep{lester2021power} learn virtual tokens that occupy context positions. Side-tuning \citep{zhang2019side} trains a parallel network fused at the output layer. All require gradient descent for each adaptation. An enmeshed network adapts at inference cost---one forward pass through $\phi$ produces new modulation tensors---and fuses at every patched layer rather than only at the output. The mechanistic similarity to LoRA (both apply low-rank perturbations) is precise but the operational difference is fundamental: LoRA is a training-time method that produces fixed weights; enmeshed modulation is an inference-time method that produces experience-specific tensors. Soft prompts and prefix tuning are the closest context-level analogue, but they still occupy context positions: virtual tokens compete for attention and are inspectable by the model's reasoning process. Compiled M operates on pre-projection hidden states, below the level of the model's chain-of-thought, providing non-inspectability that no prompt-level method can achieve (\S\ref{sec:why_not_prompt}).

\textbf{Activation steering and representation engineering.} Representation engineering \citep{zou2023representation} and activation addition \citep{turner2023activation} extract fixed directions from contrastive data and inject them during inference. These methods are elegant and effective for unidimensional behavioral shifts. However, the directions are non-adaptive: the same vector is applied regardless of context or experience. In HEXIS, the directional channel $d^*$ uses standard representation engineering for stance direction, while the enmeshed channel M provides the adaptive, multidimensional, experience-specific component. We validate empirically that the channels are orthogonal ($\Delta_\text{dir} = -0.001$ nats).

\textbf{Fast weight programmers.} \citet{schmidhuber1992learning} introduced networks that write to their own weights during a forward pass. The outer network produces weight updates for the inner network, enabling one-shot learning without backpropagation through the inner loop. HEXIS shares this spirit---$\phi$ produces modulation tensors that alter the host's computation---but differs in three ways: (1) modulation targets attention geometry (Q/V) rather than arbitrary weights, (2) the host model is a modern frozen LLM rather than a small trained-from-scratch network, and (3) compilation from structured schemas replaces online weight writing.

\textbf{Memory and identity in cognitive science.} The explicit/implicit memory distinction \citep{graf1985implicit, schacter1987implicit} provides architectural \emph{analogy} for the two-channel framing: a chain-of-thought-accessible channel for content and a perception-shaping channel for disposition. We do not claim that compiled M satisfies Schacter's neuroscience-derived criteria for implicit memory in a strict sense --- in particular, M's effect on hidden states is decodable by an external linear probe (Appendix~\ref{sec:implicit_validation}), as is true of any weight-space perturbation including a LoRA adapter. What HEXIS does provide and what motivates the analogy is \emph{architectural dissociation under adversarial pressure}: explicit beliefs in context dilute with length (\S\ref{sec:dilution}) and fold under sycophantic argument (\S\ref{sec:sycophancy}); compiled M is structurally immune to both because it lives outside the attention window the adversary can argue within. The Mind Tree's cognitive schema draws on structured self-models from developmental psychology \citep{neisser1988five} and the Toulmin model of argumentation \citep{toulmin1958uses}.
